% !TEX root = ../../main_without_quantization.tex
\section{Specialized Architectures}
\label{sec:specialized_architectures}

While Multi-Layer Perceptrons offer a generic learning model over vector inputs, they do not inherently leverage any underlying structure of the data, and thus can only indirectly learn the temporal dependencies of sequences, the latent structure of data compressed in autoencoders, or the relational structure of positions within a context in token-based models. This fundamental limitation gives rise to architectures tailored to the structure of the data.

\subsection{Recurrent Neural Networks (RNNs)}
RNNs are designed for sequential data (e.g., time series, text). Unlike feedforward networks, RNNs utilize cyclic connections to maintain a \textit{hidden state} that captures historical information. The network dynamics are described by:
\[
h_t = \sigma(W_{hh} h_{t-1} + W_{xh} x_t + b_h)
\]
\[
y_t = W_{hy} h_t + b_y
\]
where \( h_t \) is the hidden state, \( x_t \) the input, and \( y_t \) the output at time \( t \).
Figure~\ref{fig:rnn-compact} represents this recurrence by unrolling the same recurrent computation across time steps, so that the dependence of $h_t$ on $h_{t-1}$ becomes visible.
\begin{figure}[htbp]
  \centering
  \begin{tikzpicture}[>=Stealth, node distance=10mm and 18mm, every node/.style={font=\small}]
    \tikzstyle{inp} = [circle, draw=slateblue, fill=mistyblue!20, minimum size=7mm]
    \tikzstyle{hid} = [rectangle, draw=dustyteal, fill=sagegreen!20, minimum size=9mm]
    \tikzstyle{outnode} = [circle, draw=slateblue, fill=mistyblue!20, minimum size=7mm]
    \tikzstyle{lbl} = [font=\scriptsize, color=slateblue]
    
    % Inputs (bottom layer)
    \node[inp] (x1) {$x_1$};
    \node[inp] (x2) [right=of x1] {$x_2$};
    \node[inp] (x3) [right=of x2] {$x_3$};
    \node (xd) [right=of x3] {$\dots$};
    \node[inp] (xT) [right=of xd] {$x_T$};
    
    % Hidden (middle layer)
    \node[hid] (h1) [above=of x1] {$h_1$};
    \node[hid] (h2) [above=of x2] {$h_2$};
    \node[hid] (h3) [above=of x3] {$h_3$};
    \node[hid] (hT) [above=of xT] {$h_T$};
    \node (hd) [right=of h3] {$\dots$};
    
    % Outputs (top layer)
    \node[outnode] (y1) [above=of h1] {$y_1$};
    \node[outnode] (y2) [above=of h2] {$y_2$};
    \node[outnode] (y3) [above=of h3] {$y_3$};
    \node[outnode] (yT) [above=of hT] {$y_T$};
    
    % Input to hidden
    \draw[->] (x1) -- node[midway, left, lbl] {$W_{xh}$} (h1);
    \draw[->] (x2) -- node[midway, left, lbl] {$W_{xh}$} (h2);
    \draw[->] (x3) -- node[midway, left, lbl] {$W_{xh}$} (h3);
    \draw[->] (xT) -- node[midway, left, lbl] {$W_{xh}$} (hT);
    
    % Hidden recurrence 
    \draw[->] (h1) to[bend left=8] node[midway, above, lbl] {$W_{hh}$} (h2);
    \draw[->] (h2) to[bend left=8] node[midway, above, lbl] {$W_{hh}$} (h3);
    \draw[->] (h3) -- (hd);
    \draw[->] (hd) to[bend left=8] node[midway, above, lbl] {$W_{hh}$} (hT);
    
    % Hidden to Output
    \draw[->] (h1) -- node[midway, left, lbl] {$W_{hy}$} (y1);
    \draw[->] (h2) -- node[midway, left, lbl] {$W_{hy}$} (y2);
    \draw[->] (h3) -- node[midway, left, lbl] {$W_{hy}$} (y3);
    \draw[->] (hT) -- node[midway, left, lbl] {$W_{hy}$} (yT);
    
    \node[below left=2mm and 6mm of h1, font=\scriptsize] (h0label) {$h_0$};
    \draw[->] (h0label.north) to[bend left=20] (h1.west);
  \end{tikzpicture}
  \caption{Recurrent Neural Network (RNN) unrolled in time. The input $x_t$ at each time step updates the hidden state $h_t$ through the recurrent weights $W_{hh}$, while $W_{xh}$ maps inputs to hidden states and $W_{hy}$ maps hidden states to outputs $y_t$.}
  \label{fig:rnn-compact}
\end{figure}

RNNs focus on preserving information across time.
Another important architectural idea is to preserve only the most essential information by compressing the input into a lower-dimensional representation.

\subsection{Autoencoders (AEs)}

Autoencoders are neural networks designed for unsupervised learning of efficient data codings. They act as an identity function that learns to compress and then reconstruct data. Figure~\ref{fig:autoencoder} shows this structure as an encoder that maps the input to a latent space and a decoder that reconstructs the output from that compressed representation.

\begin{figure}[htbp]
    \centering
    \begin{tikzpicture}[>=Stealth, node distance=15mm, every node/.style={font=\small}]
        
        \tikzstyle{neuron} = [circle, draw=slateblue, fill=mistyblue!20, minimum size=7mm, inner sep=0pt]
        \tikzstyle{bottleneck} = [circle, draw=dustyteal, fill=sagegreen!20, minimum size=7mm, inner sep=0pt]
        \tikzstyle{lbl} = [font=\scriptsize, color=slateblue]

        % --- Nodes ---
        
        % Input Layer (x)
        \node[neuron] (i1) at (0, 2.0) {};
        \node[neuron] (i2) at (0, 1.0) {};
        \node[neuron] (i3) at (0, 0.0) {};
        \node[neuron] (i4) at (0, -1.0) {};
        \node[neuron] (i5) at (0, -2.0) {};
        
        % Encoder Hidden
        \node[neuron] (e1) at (2, 1.2) {};
        \node[neuron] (e2) at (2, 0) {};
        \node[neuron] (e3) at (2, -1.2) {};

        % Bottleneck / Latent Space (z)
        \node[bottleneck] (z1) at (4, 0.5) {};
        \node[bottleneck] (z2) at (4, -0.5) {};

        % Decoder Hidden
        \node[neuron] (d1) at (6, 1.2) {};
        \node[neuron] (d2) at (6, 0) {};
        \node[neuron] (d3) at (6, -1.2) {};

        % Output Layer (x_hat)
        \node[neuron] (o1) at (8, 2.0) {};
        \node[neuron] (o2) at (8, 1.0) {};
        \node[neuron] (o3) at (8, 0.0) {};
        \node[neuron] (o4) at (8, -1.0) {};
        \node[neuron] (o5) at (8, -2.0) {};

        % --- Connections ---
        
        % Input -> Encoder
        \foreach \i in {1,...,5}
            \foreach \j in {1,...,3}
                \draw[->] (i\i) -- (e\j);

        % Encoder -> Bottleneck
        \foreach \i in {1,...,3}
            \foreach \j in {1,...,2}
                \draw[->] (e\i) -- (z\j);
                
        % Bottleneck -> Decoder
        \foreach \i in {1,...,2}
            \foreach \j in {1,...,3}
                \draw[->] (z\i) -- (d\j);

        % Decoder -> Output
        \foreach \i in {1,...,3}
            \foreach \j in {1,...,5}
                \draw[->] (d\i) -- (o\j);

        % --- Labels ---
        
        \node[below=0.2cm of i5] {Input $x$};
        \node[below=0.2cm of o5] {Output $\hat{x}$};
        \node[below=0.2cm of z2, align=center] {Latent\\Space $z$};
        
        % Braces for Encoder/Decoder
        \draw[decorate, decoration={brace, amplitude=10pt, raise=5pt}, thick, slateblue] (i1.north) -- (z1.north) node[midway, above=20pt] {\textbf{Encoder}};
        \draw[decorate, decoration={brace, amplitude=10pt, raise=5pt}, thick, slateblue] (z1.north) -- (o1.north) node[midway, above=20pt] {\textbf{Decoder}};

    \end{tikzpicture}
    \caption[Autoencoder architecture]{Autoencoder architecture. The encoder compresses the input $x$ into a lower-dimensional latent space $z$, also called the bottleneck, and the decoder reconstructs the output $\hat{x}$ from that latent representation.}
    \label{fig:autoencoder}
\end{figure}

\begin{itemize}
    \item \textbf{Encoder:} Maps the high-dimensional input $x$ to a lower-dimensional latent representation $z$.
    \item \textbf{Decoder:} Reconstructs the output $\hat{x}$ from the latent representation $z$.
\end{itemize}

The network is trained to minimize the reconstruction error (e.g., Mean Squared Error) between the input and the output:
\[
L(x, \hat{x}) = || x - \hat{x} ||^2
\]
By forcing data through a bottleneck, the network learns to capture the most salient features of the data distribution.

Autoencoders show how architectural structure can impose a useful compression constraint.
For the language and sequence models considered later in this report, the most important specialized architecture is the Transformer.




\subsection{Transformer Networks}

Transformer networks, introduced by \cite{vaswani2017attention}, changed sequence modeling by replacing recurrent processing with attention-based computation. Instead of processing tokens one after another, a Transformer can compare tokens across the whole sequence in parallel. This makes it highly effective for language modelling, translation, summarization, and other sequence tasks. The original architecture was an encoder-decoder model for machine translation, while later systems often use encoder-only variants such as BERT \cite{devlin2019bert} or decoder-only variants such as LLaMA \cite{touvron2023llama2}.

\subsection{Core Architectural Components: Building Blocks of the Transformer}

Before separating the model into encoder and decoder stacks, it is useful to define the components repeated inside those stacks. Transformer layers are built from token embeddings, positional information, attention projections, feed-forward channels, residual connections, and normalization layers. These components are also the main targets for compression, because they determine both parameter count and runtime cost.

\subsubsection{Input Representation: Embeddings and Positional Encoding}

Every token is mapped to a continuous embedding vector. Since self-attention does not inherently encode token order, positional information is added to the token representations. The original Transformer uses sinusoidal positional encodings \cite{vaswani2017attention}:
\begin{equation}
    PE_{(\text{pos}, 2i)} = \sin\left(\frac{\text{pos}}{10000^{2i/d_{\text{model}}}}\right), \quad
    PE_{(\text{pos}, 2i+1)} = \cos\left(\frac{\text{pos}}{10000^{2i/d_{\text{model}}}}\right).
\end{equation}

\subsubsection{Self-Attention Mechanism}

Self-attention allows each token representation to aggregate information from the rest of the sequence. The mechanism uses three learned projections of each input: queries, keys, and values \cite{vaswani2017attention}. Queries represent what a token is looking for, keys represent what each token offers for comparison, and values contain the information that will be aggregated.

For an input sequence representation $\matr{X} \in \reals^{n \times d_{\text{model}}}$, the projections are:
\begin{equation}
    \matr{Q} = \matr{X}\matr{W}^{Q}, \quad
    \matr{K} = \matr{X}\matr{W}^{K}, \quad
    \matr{V} = \matr{X}\matr{W}^{V}.
\end{equation}

Scaled dot-product attention is then computed as:
\begin{equation}
    \operatorname{Attention}(\matr{Q}, \matr{K}, \matr{V})
    =
    \operatorname{softmax}\left(\frac{\matr{Q}\matr{K}^{\top}}{\sqrt{d_k}}\right)\matr{V}.
\end{equation}
The factor $\sqrt{d_k}$ prevents dot products from becoming too large, which would push the softmax into saturated regions with weak gradients.

\subsubsection{Multi-Head Attention}

Multi-head attention applies several attention mechanisms in parallel, allowing the model to attend to different representation subspaces and token relationships at the same time \cite{vaswani2017attention}. For $h$ heads:
\begin{equation}
    \operatorname{head}_i =
    \operatorname{Attention}(\matr{X}\matr{W}^{Q}_{i}, \matr{X}\matr{W}^{K}_{i}, \matr{X}\matr{W}^{V}_{i}),
    \quad i = 1,\ldots,h.
\end{equation}

The head outputs are concatenated and projected back to the model dimension:
\begin{equation}
    \operatorname{MultiHead}(\matr{X}) =
    \operatorname{Concat}(\operatorname{head}_1,\ldots,\operatorname{head}_h)\matr{W}^{O}.
\end{equation}
Here, $\matr{W}^{O} \in \reals^{hd_v \times d_{\text{model}}}$ is the output projection. Attention heads and their projection matrices are important pruning targets because entire heads can sometimes be removed with limited accuracy loss.

\subsubsection{Feed-Forward Network}

After attention mixes information across tokens, a position-wise Feed-Forward Network (FFN) processes each token representation independently:
\begin{equation}
    \operatorname{FFN}(\vect{x}) = \matr{W}_2 \sigma(\matr{W}_1\vect{x} + \vect{b}_1) + \vect{b}_2.
\end{equation}
The FFN usually expands the hidden representation to a larger intermediate dimension $d_{ff}$ and then projects it back to $d_{\text{model}}$. In the original Transformer, $d_{ff}$ is typically much larger than $d_{\text{model}}$, making FFN channels a major source of parameters and a natural target for structured pruning.

\subsubsection{Residual Connections and Layer Normalization}

Each attention or FFN sub-layer is wrapped by a residual path and layer normalization. Residual connections preserve a direct path for information and gradients, while normalization stabilizes the scale of hidden representations. This design is essential for training deep Transformer stacks and is also important when later compression methods remove heads, channels, or blocks \cite{vaswani2017attention,He2016,Ba2016}.

\subsection{Encoder and Decoder Architecture}

With the core components defined, the encoder and decoder can be described cleanly. The full encoder-decoder Transformer contains two main stacks:
\begin{itemize}
    \item \textbf{Encoder:} Processes the complete input sequence in parallel and produces contextual source representations.
    \item \textbf{Decoder:} Generates the output sequence autoregressively while attending to both previous target tokens and the encoder output.
\end{itemize}

The encoder converts the source sequence into a contextual memory representation $\matr{H}_{\text{enc}}$. The decoder uses this memory through cross-attention while also using masked self-attention to prevent access to future target tokens. The complete architecture is shown in Figure~\ref{fig:transformer-architecture}.

\begin{figure}[H]
  \centering
  \begin{tikzpicture}[scale=1.2, transform shape]
    \tikzstyle{Sanode} = [minimum height=1.4em, minimum width=7em, inner sep=3pt, rounded corners=1.5pt, draw=slateblue, fill=mistyblue!30]
    \tikzstyle{Resnode} = [minimum height=1.1em, minimum width=7em, inner sep=3pt, rounded corners=1.5pt, draw=dustyteal, fill=softgray!30]
    \tikzstyle{ffnnode} = [minimum height=1.4em, minimum width=7em, inner sep=3pt, rounded corners=1.5pt, draw=slateblue, fill=sagegreen!30]
    \tikzstyle{outputnode} = [minimum height=1.4em, minimum width=7em, inner sep=3pt, rounded corners=1.5pt, draw=slateblue, fill=mistyblue!50]
    \tikzstyle{inputnode} = [minimum height=1.4em, minimum width=3.5em, inner sep=3pt, rounded corners=1.5pt, draw=dustyteal, fill=mistyblue!20]
    \tikzstyle{posnode} = [minimum height=1.4em, minimum width=3.5em, inner sep=3pt, rounded corners=1.5pt, draw=dustyteal, fill=softgray!50]
    \tikzstyle{standard} = [rounded corners=3pt]

    \node [Sanode, anchor=west] (sa1) at (0,0) {\tiny \textbf{Self-Attention}};
    \node [Resnode, anchor=south] (res1) at ([yshift=0.3em]sa1.north) {\tiny \textbf{Add \& LayerNorm}};
    \node [ffnnode, anchor=south] (ffn1) at ([yshift=1em]res1.north) {\tiny \textbf{FFN}};
    \node [Resnode, anchor=south] (res2) at ([yshift=0.3em]ffn1.north) {\tiny \textbf{Add \& LayerNorm}};
    \node [inputnode, anchor=north west] (input1) at ([yshift=-1em,xshift=-0.5em]sa1.south west) {\tiny \textbf{Emb.}};
    \node [] (add1) at ([yshift=-1.6em,xshift=3.5em]sa1.south west) {$+$};
    \node [posnode, anchor=north east] (pos1) at ([yshift=-1em,xshift=0.5em]sa1.south east) {\tiny \textbf{Position}};
    \node [anchor=south] (encoder) at ([xshift=0.2em,yshift=0.6em]res2.north west) {\scriptsize \textbf{Encoder}};
    \draw [->,dustyteal] ([yshift=-4em]sa1.south) -- (add1);

    \draw [->,dustyteal] (sa1.north) -- (res1.south);
    \draw [->,dustyteal] (res1.north) -- (ffn1.south);
    \draw [->,dustyteal] (ffn1.north) -- (res2.south);
    \draw [->,dustyteal] ([yshift=-1em]sa1.south) -- (sa1.south);

    \node [Sanode, anchor=west] (sa2) at ([xshift=3em]sa1.east) {\tiny \textbf{Masked Self-Attn}};
    \node [Resnode, anchor=south] (res3) at ([yshift=0.3em]sa2.north) {\tiny \textbf{Add \& LayerNorm}};
    \node [Sanode, anchor=south] (ed1) at ([yshift=1em]res3.north) {\tiny \textbf{Cross-Attention}};
    \node [Resnode, anchor=south] (res4) at ([yshift=0.3em]ed1.north) {\tiny \textbf{Add \& LayerNorm}};
    \node [ffnnode, anchor=south] (ffn2) at ([yshift=1em]res4.north) {\tiny \textbf{FFN}};
    \node [Resnode, anchor=south] (res5) at ([yshift=0.3em]ffn2.north) {\tiny \textbf{Add \& LayerNorm}};
    \node [outputnode, anchor=south] (o1) at ([yshift=1em]res5.north) {\tiny \textbf{Linear + Softmax}};
    \node [inputnode, anchor=north west] (input2) at ([yshift=-1em,xshift=-0.5em]sa2.south west) {\tiny \textbf{Emb.}};
    \node [] (add) at ([yshift=-1.6em,xshift=3.5em]sa2.south west) {$+$};
    \node [posnode, anchor=north east] (pos2) at ([yshift=-1em,xshift=0.5em]sa2.south east) {\tiny \textbf{Position}};
    \node [anchor=east] (decoder) at ([xshift=-1em,yshift=-1.5em]o1.west) {\scriptsize \textbf{Decoder}};

    \draw[<-, dustyteal] (add.south) -- ++(0,-1.5em);

    \draw [->,dustyteal] (sa2.north) -- (res3.south);
    \draw [->,dustyteal] (res3.north) -- (ed1.south);
    \draw [->,dustyteal] (ed1.north) -- (res4.south);
    \draw [->,dustyteal] (res4.north) -- (ffn2.south);
    \draw [->,dustyteal] (ffn2.north) -- (res5.south);
    \draw [->,dustyteal] (res5.north) -- (o1.south);
    \draw [->,dustyteal] (o1.north) -- ([yshift=0.5em]o1.north);
    \draw [->,dustyteal] ([yshift=-1em]sa2.south) -- (sa2.south);

    \draw[->,standard,dustyteal] ([yshift=-0.5em]sa1.south) -- ([xshift=-4em,yshift=-0.5em]sa1.south) -- ([xshift=-4em,yshift=2.3em]sa1.south) -- ([xshift=-3.5em,yshift=2.3em]sa1.south);
    \draw[->,standard,dustyteal] ([yshift=0.5em]res1.north) -- ([xshift=-4em,yshift=0.5em]res1.north) -- ([xshift=-4em,yshift=3.3em]res1.north) -- ([xshift=-3.5em,yshift=3.3em]res1.north);
    \draw[->,standard,dustyteal] ([yshift=-0.5em]sa2.south) -- ([xshift=4em,yshift=-0.5em]sa2.south) -- ([xshift=4em,yshift=2.3em]sa2.south) -- ([xshift=3.5em,yshift=2.3em]sa2.south);
    \draw[->,standard,dustyteal] ([yshift=0.5em]res3.north) -- ([xshift=4em,yshift=0.5em]res3.north) -- ([xshift=4em,yshift=3.3em]res3.north) -- ([xshift=3.5em,yshift=3.3em]res3.north);
    \draw[->,standard,dustyteal] ([yshift=0.5em]res4.north) -- ([xshift=4em,yshift=0.5em]res4.north) -- ([xshift=4em,yshift=3.3em]res4.north) -- ([xshift=3.5em,yshift=3.3em]res4.north);
    \draw[->,standard, line width=1.2pt, slateblue] (res2.north) -- ([yshift=0.5em]res2.north) -- ([xshift=5em,yshift=0.5em]res2.north) -- ([xshift=5em,yshift=-2.2em]res2.north) -- ([xshift=6.5em,yshift=-2.2em]res2.north);
    \node[font=\tiny, slateblue] at ([xshift=5.5em,yshift=-1.5em]res2.north) {$\matr{H}_{\text{enc}}$};

  \end{tikzpicture}
  \caption{Complete Encoder-Decoder Transformer architecture. The encoder memory matrix $\matr{H}_{\text{enc}}$ (highlighted in blue) is fed to all decoder layers for cross-attention. Residual connections (curved arrows) enable stable gradient flow throughout the network.}
  \label{fig:transformer-architecture}
\end{figure}

\subsubsection{The Encoder Stack}

An encoder layer contains multi-head self-attention followed by a feed-forward network, with residual connections and normalization around each sub-layer. In encoder self-attention, the queries, keys, and values all come from the same source sequence. This makes the encoder bidirectional: every input token can attend to every other input token. Figure~\ref{fig:encoder} shows the internal order of these operations before the resulting representation is passed to the decoder. After $N$ encoder layers, the output is the memory matrix $\matr{H}_{\text{enc}} \in \reals^{n \times d_{\text{model}}}$, which summarizes the source sequence for the decoder.

\begin{figure}[ht]
  \centering
  \begin{tikzpicture}[scale=1.2, transform shape]
    \tikzstyle{Sanode} = [minimum height=1.4em, minimum width=7em, inner sep=3pt, rounded corners=1.5pt, draw=slateblue, fill=mistyblue!30]
    \tikzstyle{Resnode} = [minimum height=1.1em, minimum width=7em, inner sep=3pt, rounded corners=1.5pt, draw=dustyteal, fill=softgray!30]
    \tikzstyle{ffnnode} = [minimum height=1.4em, minimum width=7em, inner sep=3pt, rounded corners=1.5pt, draw=slateblue, fill=sagegreen!30]
    \tikzstyle{inputnode} = [minimum height=1.4em, minimum width=3.5em, inner sep=3pt, rounded corners=1.5pt, draw=dustyteal, fill=mistyblue!20]
    \tikzstyle{posnode} = [minimum height=1.4em, minimum width=3.5em, inner sep=3pt, rounded corners=1.5pt, draw=dustyteal, fill=softgray!50]
    \tikzstyle{standard} = [rounded corners=3pt]

    \node [Sanode, anchor=west] (sa1) at (0,0) {\tiny \textbf{Self-Attention}};
    \node [Resnode, anchor=south] (res1) at ([yshift=0.3em]sa1.north) {\tiny \textbf{Add \& LayerNorm}};
    \node [ffnnode, anchor=south] (ffn1) at ([yshift=1em]res1.north) {\tiny \textbf{FFN}};
    \node [Resnode, anchor=south] (res2) at ([yshift=0.3em]ffn1.north) {\tiny \textbf{Add \& LayerNorm}};
    \node [inputnode, anchor=north west] (input1) at ([yshift=-1em,xshift=-0.5em]sa1.south west) {\tiny \textbf{Emb.}};
    \node [] (add1) at ([yshift=-1.6em,xshift=3.5em]sa1.south west) {$+$};
    \node [posnode, anchor=north east] (pos1) at ([yshift=-1em,xshift=0.5em]sa1.south east) {\tiny \textbf{Position}};
    \node [anchor=south] (encoder) at ([xshift=0.2em,yshift=0.6em]res2.north west) {\scriptsize \textbf{Encoder}};

    \draw[<-, dustyteal] (add1.south) -- ++(0,-1.5em);

    \draw [->,dustyteal] (sa1.north) -- (res1.south);
    \draw [->,dustyteal] (res1.north) -- (ffn1.south);
    \draw [->,dustyteal] (ffn1.north) -- (res2.south);
    \draw [->,dustyteal] ([yshift=-1em]sa1.south) -- (sa1.south);

    \draw[->,standard,dustyteal] ([yshift=-0.5em]sa1.south) -- ([xshift=-4em,yshift=-0.5em]sa1.south) -- ([xshift=-4em,yshift=2.3em]sa1.south) -- ([xshift=-3.5em,yshift=2.3em]sa1.south);
    \draw[->,standard,dustyteal] ([yshift=0.5em]res1.north) -- ([xshift=-4em,yshift=0.5em]res1.north) -- ([xshift=-4em,yshift=3.3em]res1.north) -- ([xshift=-3.5em,yshift=3.3em]res1.north);

  \end{tikzpicture}
  \caption{Transformer encoder layer architecture. The input embedding (Emb.) is combined with positional information, processed by self-attention and a Feed-Forward Network (FFN), and stabilized by residual connections with Layer Normalization (LayerNorm).}
  \label{fig:encoder}
\end{figure}

\subsubsection{The Decoder Stack}

A decoder layer adds masked self-attention and cross-attention before the feed-forward network, enabling autoregressive generation while attending to encoder outputs. Masked self-attention restricts token $t$ to positions $\leq t$, preventing information leakage from future target tokens. Cross-attention then links the decoder to the encoder: decoder states act as queries, while encoder states provide keys and values. Figure~\ref{fig:decoder} shows this additional cross-attention block, which is the main structural difference between the decoder layer and the encoder layer.

\begin{figure}[htbp]
  \centering
  \begin{tikzpicture}[scale=1.2, transform shape]
    \tikzstyle{Sanode} = [minimum height=1.4em, minimum width=7em, inner sep=3pt, rounded corners=1.5pt, draw=slateblue, fill=mistyblue!30]
    \tikzstyle{Resnode} = [minimum height=1.1em, minimum width=7em, inner sep=3pt, rounded corners=1.5pt, draw=dustyteal, fill=softgray!30]
    \tikzstyle{ffnnode} = [minimum height=1.4em, minimum width=7em, inner sep=3pt, rounded corners=1.5pt, draw=slateblue, fill=sagegreen!30]
    \tikzstyle{outputnode} = [minimum height=1.4em, minimum width=7em, inner sep=3pt, rounded corners=1.5pt, draw=slateblue, fill=mistyblue!50]
    \tikzstyle{inputnode} = [minimum height=1.4em, minimum width=3.5em, inner sep=3pt, rounded corners=1.5pt, draw=dustyteal, fill=mistyblue!20]
    \tikzstyle{posnode} = [minimum height=1.4em, minimum width=3.5em, inner sep=3pt, rounded corners=1.5pt, draw=dustyteal, fill=softgray!50]
    \tikzstyle{standard} = [rounded corners=3pt]

    \node [Sanode, anchor=west] (sa2) at (0,0) {\tiny \textbf{Masked Self-Attn}};
    \node [Resnode, anchor=south] (res3) at ([yshift=0.3em]sa2.north) {\tiny \textbf{Add \& LayerNorm}};
    \node [Sanode, anchor=south] (ed1) at ([yshift=1em]res3.north) {\tiny \textbf{Cross-Attention}};
    \node [Resnode, anchor=south] (res4) at ([yshift=0.3em]ed1.north) {\tiny \textbf{Add \& LayerNorm}};
    \node [ffnnode, anchor=south] (ffn2) at ([yshift=1em]res4.north) {\tiny \textbf{FFN}};
    \node [Resnode, anchor=south] (res5) at ([yshift=0.3em]ffn2.north) {\tiny \textbf{Add \& LayerNorm}};
    \node [outputnode, anchor=south] (o1) at ([yshift=1em]res5.north) {\tiny \textbf{Linear + Softmax}};
    \node [inputnode, anchor=north west] (input2) at ([yshift=-1em,xshift=-0.5em]sa2.south west) {\tiny \textbf{Emb.}};
    \node [] (add) at ([yshift=-1.6em,xshift=3.5em]sa2.south west) {$+$};
    \node [posnode, anchor=north east] (pos2) at ([yshift=-1em,xshift=0.5em]sa2.south east) {\tiny \textbf{Position}};
    \node [anchor=south] (decoder) at ([xshift=0.2em,yshift=0.6em]o1.north west) {\scriptsize \textbf{Decoder}};

    \draw [->,dustyteal] ([yshift=-4em]sa2.south) -- (add);

    \draw [->,dustyteal] (sa2.north) -- (res3.south);
    \draw [->,dustyteal] (res3.north) -- (ed1.south);
    \draw [->,dustyteal] (ed1.north) -- (res4.south);
    \draw [->,dustyteal] (res4.north) -- (ffn2.south);
    \draw [->,dustyteal] (ffn2.north) -- (res5.south);
    \draw [->,dustyteal] (res5.north) -- (o1.south);
    \draw [->,dustyteal] (o1.north) -- ([yshift=0.5em]o1.north);
    \draw [->,dustyteal] ([yshift=-1em]sa2.south) -- (sa2.south);

    \draw[->,standard,dustyteal] ([yshift=-0.5em]sa2.south) -- ([xshift=4em,yshift=-0.5em]sa2.south) -- ([xshift=4em,yshift=2.3em]sa2.south) -- ([xshift=3.5em,yshift=2.3em]sa2.south);
    \draw[->,standard,dustyteal] ([yshift=0.5em]res3.north) -- ([xshift=4em,yshift=0.5em]res3.north) -- ([xshift=4em,yshift=3.3em]res3.north) -- ([xshift=3.5em,yshift=3.3em]res3.north);
    \draw[->,standard,dustyteal] ([yshift=0.5em]res4.north) -- ([xshift=4em,yshift=0.5em]res4.north) -- ([xshift=4em,yshift=3.3em]res4.north) -- ([xshift=3.5em,yshift=3.3em]res4.north);

  \end{tikzpicture}
  \caption{Transformer decoder layer architecture. The decoder contains masked self-attention for previous target tokens, cross-attention for reading encoder outputs, and a Feed-Forward Network (FFN), each wrapped with residual connections and Layer Normalization (LayerNorm).}
  \label{fig:decoder}
\end{figure}

\subsubsection{Encoder-Decoder Information Flow}

The complete sequence-to-sequence computation proceeds as follows:
\begin{enumerate}[leftmargin=*]
  \item \textbf{Encoding phase.} Source tokens $\vect{x}=(x_1,\ldots,x_n)$ are embedded, combined with positional information, and passed through $N$ encoder layers. The output is $\matr{H}_{\text{enc}}$.
  \item \textbf{Decoding phase.} Target tokens $\vect{y}=(y_1,\ldots,y_t)$ are embedded and processed autoregressively. Each decoder layer applies masked self-attention, cross-attention over $\matr{H}_{\text{enc}}$, and a feed-forward transformation.
  \item \textbf{Output phase.} The final decoder representation is projected to vocabulary logits. A softmax layer converts logits into token probabilities, and the selected token is fed back into the decoder at the next step.
\end{enumerate}

\vspace{0.8em}
\noindent\rule{\textwidth}{0.4pt}
\vspace{0.8em}

At this point in the report, we can view the Transformer as an arrangement of embedding layers, attention layers, feed-forward functions, and encoder--decoder connections. This architecture accounts for why it performs well on sequence modelling tasks, since each position can refine its representation using information from the rest of the sequence. This expressive structure, however, is computationally demanding. The next point examines where this cost arises inside the Transformer architecture.

\subsection{The Computational Bottleneck of Transformers}
\label{sec:transformer_bottleneck}
While powerful, the Transformer's success comes at a significant computational cost, which forms a central motivation for the compression techniques explored in this report. The standard Transformer architecture replaces recurrence with self-attention, enabling strong parallelism and high performance across sequence modelling tasks~\cite{vaswani2017attention}. However, this advantage comes with a major bottleneck: in full self-attention, each token attends to every other token in the sequence. For a sequence of length $n$, this requires an attention score matrix of size $n \times n$, making the attention component scale quadratically with sequence length in both computation and memory~\cite{beltagy2020longformer,tay2020efficient}.

This quadratic behaviour has direct practical implications. Doubling the sequence length from $512$ to $1024$ multiplies the attention matrix size by four, since $1024^2 / 512^2 = 4$. As a result, processing long documents or long-context inputs becomes substantially more expensive, which has motivated a large body of work on efficient Transformer variants and sparse or linear attention mechanisms~\cite{beltagy2020longformer,tay2020efficient}.

At the same time, modern language models benefit strongly from scaling model size, training data, and compute~\cite{kaplan2020scaling}. Later scaling-law work further showed that model size and the number of training tokens should be scaled together under a fixed compute budget, highlighting the continuing importance of large-scale training~\cite{hoffmann2022training}. This creates a tension between scale and deployment: larger models often achieve better performance, but practical inference is constrained by latency, memory, energy, and hardware availability. This tension motivates compression methods such as pruning, quantization, and knowledge distillation~\cite{han2015deep,sanh2019distilbert}.

\par\smallskip
\noindent\rule{0.35\linewidth}{0.4pt}
\par\smallskip

The computational bottleneck defined above can be considered in several different ways. A Transformer may be expensive because of the size of its parameters, the memory required for activations, the number of arithmetic operations it performs, or the time required to run it on a specific hardware platform. Before discussing model efficiency in more detail, it is therefore useful to introduce the most common measures used to assess computational cost.

\subsection{Hardware-Aware Efficiency Metrics}

Compression is evaluated through several distinct efficiency metrics. Parameter count, memory footprint, arithmetic cost, latency, and throughput capture different aspects of model execution. Hardware-aware evaluation therefore requires separate treatment of storage cost, arithmetic cost, and runtime behavior~\cite{Sze2017EfficientDNN}.

\subsubsection{Parameter Count and Model Size}

Let $\params = \{W^{(l)}, b^{(l)}\}_{l=1}^{L}$ denote the parameters of a network. The total parameter count is
\begin{equation}
P = \sum_{l=1}^{L} \left( |W^{(l)}| + |b^{(l)}| \right).
\end{equation}
If the parameters are stored with precision $b_w$ bits, the static model footprint is
\begin{equation}
M_{\text{model}} = \frac{P\, b_w}{8} \quad \text{bytes}.
\end{equation}
This follows from counting the stored elements and multiplying by the number of bytes used by each element~\cite{PyTorchNumel,PyTorchElementSize}.
This quantity determines whether a model fits within device memory and how much parameter traffic must be sustained during inference. Quantization reduces $b_w$, pruning reduces the number of active parameters, and low-rank factorization reduces the parameterization itself.

For sparse models, storage depends on both the remaining weights and the associated metadata. If $P_{\text{nz}}$ denotes the number of non-zero weights and $M_{\text{meta}}$ the storage needed for indices or masks, then
\begin{equation}
M_{\text{sparse}} = \frac{P_{\text{nz}}\, b_w}{8} + M_{\text{meta}}.
\end{equation}

\subsubsection{MACs, FLOPs, and Arithmetic Intensity}

For a dense matrix multiplication $Y = AB$ with $A \in \reals^{m \times k}$ and $B \in \reals^{k \times n}$, the arithmetic cost is
\begin{equation}
\mathrm{MACs}(A,B) = mkn,
\qquad
\mathrm{FLOPs}(A,B) \approx 2mkn.
\end{equation}
The distinction is conventional: one multiply-accumulate corresponds to one MAC, while FLOP counting often treats the multiply and the addition as two floating-point operations~\cite{NVIDIAMatrixMultiplication}. In Transformer layers, the feed-forward subnetwork contributes large dense matrix products, while self-attention adds the sequence-dependent $O(n^2 d_{\text{model}})$ cost of score computation and value aggregation.

Arithmetic cost alone is insufficient for hardware analysis. A useful complementary quantity is the arithmetic intensity
\begin{equation}
I = \frac{\mathrm{FLOPs}}{\mathrm{Bytes\ Moved}},
\end{equation}
which measures how much computation is performed per byte transferred~\cite{Williams2009Roofline,NVIDIAMatrixMultiplication}. Low arithmetic intensity indicates a memory-bound workload, whereas high arithmetic intensity indicates a compute-bound workload.

\subsubsection{Bandwidth, Latency, and Throughput}

Memory bandwidth measures the rate at which data can be transferred between memory and compute units. If $B_{\text{move}}$ bytes are transferred during an operation that takes time $T$, the effective bandwidth demand is
\begin{equation}
\mathrm{BW}_{\text{eff}} = \frac{B_{\text{move}}}{T}.
\end{equation}
This follows the standard effective-bandwidth definition used for GPU kernels~\cite{CUDABestPractices}.
When parameter streaming and activation movement dominate execution, bandwidth becomes the primary runtime constraint.

Latency and throughput characterize runtime from complementary viewpoints. Latency is the time required to complete one inference request or one decoding step~\cite{Sze2017EfficientDNN}:
\begin{equation}
T_{\text{latency}} = T_{\text{compute}} + T_{\text{memory}} + T_{\text{overhead}}.
\end{equation}
Throughput measures completed work per unit time, for example requests per second or generated tokens per second~\cite{Sze2017EfficientDNN}:
\begin{equation}
\mathrm{Throughput} = \frac{N_{\text{outputs}}}{T}.
\end{equation}
In large autoregressive models, latency and throughput often follow different scaling trends. Batching can increase throughput while single-request latency remains high. Parameter compression can reduce memory transfer while leaving runtime gains limited when execution is poorly aligned with the hardware schedule.
